\documentclass[11pt]{article}

% ========== LANGUAGE AND ENCODING ==========
\usepackage[utf8]{inputenc}      % Allows UTF-8 input (e.g., ć, č, š)
\usepackage[T1]{fontenc}         % Correct output of special characters
\usepackage[serbian]{babel}      % Serbian language (Latin script by default)

% ========== MATH & UTILITIES ==========
\usepackage{amsmath,amssymb}
\usepackage{xcolor}      % only load xcolor once
\usepackage{xparse}
\usepackage{enumitem}
\usepackage{tcolorbox}
\tcbuselibrary{skins,breakable}
\usepackage{titlesec}

\usepackage{pgf,tikz}
\usetikzlibrary{shadows}
\usepackage{tikz-3dplot}
\usepackage{pgfplots}
\tdplotsetmaincoords{70}{120}
\pgfplotsset{compat=1.18} % Ensures modern features are used

\titleformat{\section}[block]
  {\centering\normalfont\Huge\bfseries}
  {\thesection.}{1em}{}
\usepackage{makecell}
\usepackage{amsthm}
\usetikzlibrary{math}
\usepackage{enumitem}
\usetikzlibrary{arrows}



\pgfplotsset{compat=1.18}

\newtcolorbox{defbox}{
    colback = blue!10!white,
    colframe = blue,
    boxrule = 0.5pt,
    arc = 4pt,
    boxsep = 3pt,
    enhanced,
    drop shadow
}

\newtcolorbox{teoremabox}{
    colback = red!30!white,
    colframe = red,
    boxrule = 0.5pt,
    arc = 4pt,
    boxsep = 3pt,
    enhanced,
    drop shadow
}

\newtcolorbox{stavbox}{
    colback = orange!15!white,
    colframe = orange,
    boxrule = 0.5pt,
    arc = 4pt,
    boxsep = 3pt,
    enhanced,
    drop shadow
}

\newtcolorbox{lemabox}{
    colback = yellow!30!white,
    colframe = yellow,
    boxrule = 0.5pt,
    arc = 4pt,
    boxsep = 3pt,
    enhanced,
    drop shadow
}

\newtcolorbox{tvrbox}{
    colback = green!10!white,
    colframe = green,
    boxrule = 0.5pt,
    arc = 4pt,
    boxsep = 3pt,
    enhanced,
    drop shadow
}

\newtcolorbox{primbox}{
    colback = lime!10!white,
    colframe = lime!70!black,
    boxrule = 0.5pt,
    arc = 4pt,
    boxsep = 3pt,
    enhanced,
    drop shadow
}

\newtcolorbox{pitbox}{
    colback = orange!20!white,
    colframe = brown!70!black,
    boxrule = 0.5pt,
    arc = 4pt,
    boxsep = 3pt,
    enhanced,
    drop shadow
}

\newcommand{\real}{\mathbb{R}}

\newcommand{\gr}[2]{%
  \mathchoice
    {\bigg|_{#1}^{#2}}% display style: \[ \]
    {\big|_{#1}^{#2}}% text style: \( \)
    {\scriptstyle|_{#1}^{#2}}% script style
    {\scriptscriptstyle|_{#1}^{#2}}% scriptscript style
}
\newcommand{\sekvenca}[6]{%
\begin{tikzpicture}[scale=1.1]
    % X-axis
    \pgfmathparse{#3 + 1} \let\xend\pgfmathresult
    \draw[->] (0,0) -- (\xend,0) node[right] {#4};
    % Y-axis
    \draw[->] (0,0) -- (0,5) node[above] {$a_{#4}$};

    % Y ticks
    \foreach \y in {0.5,1,1.5,2,2.5,3,3.5,4,4.5} {
        \draw (-0.1,\y) -- (0.1,\y) node[left=6pt] {\small \y};
    }

    % Sequence points
    \foreach \n in {#2,...,#3} {
        \pgfmathparse{#1}  % expression in \n
        \let\an\pgfmathresult
        \filldraw[blue] (\n,\an) circle (2pt);
    }

    % Lower horizontal line (if provided)
    \if\relax\detokenize{#5}\relax
        % nothing
    \else
      \draw[dashed, red] (0,#5) -- (\xend,#5) node[right] {\huge ${\text{ograničenje}} = #5$};
    \fi

    % Upper horizontal line (if provided)
    \if\relax\detokenize{#6}\relax
        % nothing
    \else
      \ifx#6+infty
        % do nothing if +infty
      \else
        \draw[dashed, red] (0,#6) -- (\xend,#6) node[right] {\huge ${\text{ograničenje}} = #6$};
      \fi
    \fi
\end{tikzpicture}
}




\newcounter{mathitem} 
\renewcommand{\themathitem}{\arabic{mathitem}} 


\theoremstyle{definition}
\newtheorem{definicija}{Definicija}[section]
\newtheorem{primer}[definicija]{Primer}
\newtheorem{teorema}[definicija]{Teorema}
\newtheorem{tvrdjenje}[definicija]{Tvrđenje}
\newtheorem{lema}[definicija]{Lema}
\newtheorem{stav}[definicija]{Stav}

\makeatletter
\newtheoremstyle{serbnodot}
  {3pt}
  {3pt}
  {\normalfont}
  {}
  {\bfseries}
  {}
  {0.5em}
  {\thmname{#1}}
\makeatother

\theoremstyle{serbnodot}
\newtheorem*{dokaz}{Dokaz}


\renewcommand{\thesubsection}{\arabic{subsection}}


\titleformat{\section}[block]
  {\centering\normalfont\Huge\bfseries}
  {} 
  {0pt}
  {}


\titleformat{\subsection}[block]
  {\centering\normalfont\Large\bfseries} 
  {\thesubsection} 
  {1em} 
  {}


\NewDocumentCommand{\naslov}{O{blue!70!black} m}{%
  \begin{center}
    \vspace{1em}
    {\color{#1}\Large\bfseries #2}
    \vspace{1em}
    \par
  \end{center}
}

\newcommand{\podnaslov}[1]{%
  \vspace{2em}
  {\color{black}\Large\bfseries #1}
  \vspace{0.5em}
  \par
}

% Zadatak counter 
\newcounter{zadatak}[section]
\renewcommand{\thezadatak}{\arabic{zadatak}}

\newcommand{\pitanje}[1]{%
    \begin{pitbox}
    \textbf{Pitanje:} #1
    \end{pitbox}
}
\newcommand{\napomena}[1]{%
    \textbf{Napomena:} #1
}
% Boxed Definicija
\let\olddefinicija\definicija
\RenewDocumentEnvironment{definicija}{O{}}
{
    \begin{defbox}
    \olddefinicija[#1]
}
{
    \end{defbox}
}

% Boxed Primer
\let\oldprimer\primer
\RenewDocumentEnvironment{primer}{O{}}
{
    \begin{primbox}
    \oldprimer[#1]
}
{
    \end{primbox}
}

% Boxed Teorema
\let\oldteorema\teorema
\RenewDocumentEnvironment{teorema}{O{}}
{
    \begin{teoremabox}
    \oldteorema[#1]
}
{
    \end{teoremabox}
}

% Boxed Tvrdjenje
\let\oldtvrdjenje\tvrdjenje
\let\endoldtvrdjenje\endtvrdjenje

\RenewDocumentEnvironment{tvrdjenje}{O{}}
{
    \begin{tvrbox}
    \oldtvrdjenje[#1]\par\noindent\ignorespaces
}
{
    \endoldtvrdjenje
    \end{tvrbox}
}
% Boxed Lemma
\let\oldlema\lema
\RenewDocumentEnvironment{lema}{O{}}
{
    \begin{lemabox}
    \oldlema[#1]
}
{
    \end{lemabox}
}

% Boxed stav
\let\oldstav\stav
\RenewDocumentEnvironment{stav}{O{}}
{
    \begin{stavbox}
    \oldstav[#1]
}
{
    \end{stavbox}
}


\NewDocumentEnvironment{zadatak}{O{} O{}}
{%
  \refstepcounter{zadatak}
  \begin{tcolorbox}[
    colback=blue!3,
    colframe=blue!40!black,
    sharp corners,
    boxrule=0.8pt,
    enhanced,
    breakable,
    title={\textbf{Zadatak \thezadatak\IfValueT{#1}{. #1}}},
    #2
  ]
}
{%
  \end{tcolorbox}
}



\NewDocumentEnvironment{zadatakp}{m O{}}
{%
  \refstepcounter{zadatak}
  \begin{tcolorbox}[
    colback=blue!3,
    colframe=blue!40!black,
    sharp corners,
    boxrule=0.8pt,
    enhanced,
    breakable,
    title={\textbf{Zadatak \thezadatak. #1}},
    #2
  ]
  \begin{enumerate}[label=\alph*)]
}
{%
  \end{enumerate}
  \end{tcolorbox}
}


\usepackage[a4paper, margin=2.5cm]{geometry}


\begin{document}

\begin{titlepage}
    \centering
    \vspace*{3cm}

    {\Huge \textbf{Analiza 3}}

    \vspace{1cm}

    \rule{\textwidth}{1.2pt}

    \vspace{0.6cm}
    {\Huge \textbf{Beleške sa predavanja profesorke}}\\[0.3cm]
    {\Huge \textbf{Aleksandre Marinković}}
    \vspace{0.6cm}

    \rule{\textwidth}{1.2pt}

    \vfill

    {\large \today}

\end{titlepage}

\newpage
\tableofcontents 
\newpage

\section{Funkcije više promenljivih}
Do sada smo razmatrali nizove realnih brojeva i funkcije jedne promenljive,
zajedno sa pojmovima kao što su limes, neprekidnost i diferencijabilnost.
U ovom poglavlju ove pojmove uopštavamo na nizove u \( \mathbb{R}^n \)
i na funkcije više promenljivih.

\subsection{Prostor \( \mathbb{R}^n \)}

\begin{definicija}

    Definišemo skup \( \mathbb{R}^n \) kao skup svih \(n\)-torki realnih brojeva
    \[
        \mathbb{R}^n := \{(x_1, \dots , x_n) \mid x_i \in \mathbb{R}, i = 1, \dots , n \}.
    \]

\end{definicija}

\begin{definicija}
    Uvodimo \textbf{vektorski prostor} kao \( (\mathbb{R}^n, +, \cdot) \) gde
    elemente ovog vektorskog prostora nazivamo \textit{vektorima} koje ćemo označavati 
    \[
        x \in \mathbb{R}^n := (x_1, \dots, x_n).
    \]
    i definišemo:
    \begin{itemize}
        \item Sabiranje \( + : \mathbb{R}^n + \mathbb{R}^n \rightarrow \mathbb{R} \)

            \( x, y \in \mathbb{R}^n, \quad x + y = (x_1, \dots, x_n) + (y_1, \dots, y_n) 
            = (x_1 + y_1, \dots, x_n + y_n) \)
        \item Množenje skalarom(realnim brojem) \( \cdot : \mathbb{R} \times \mathbb{R}^n 
        \rightarrow \mathbb{R}^n \)

            \( \lambda \in \mathbb{R}, x \in \mathbb{R}^n, \quad \lambda \cdot x = 
            \lambda \cdot (x_1, \dots, x_n) = (\lambda x_1, \dots, \lambda x_n) \)
    \end{itemize}
\end{definicija}

\pitanje{
    Da li možemo izvesti oduzimanje dva vektora \( x, y \in \mathbb{R}^n \), gde je 
    \( x = (x_1, \dots, x_n) \)
    a \( y= (y_1, \dots, y_n) \)?

    \medskip
    \textbf{Da! } \( x - y = x + (-1) \cdot y\)
    \begin{align*}
        (x_1, \dots, x_n) - (y_1, \dots, y_n) &= (x_1, \dots, x_n) 
        + (-1) \cdot (y_1, \dots, y_n) \\
        &= (x_1, \dots, x_n) + (-y_1, \dots, -y_n) \\
        &= (x_1 - y_1, \dots, x_n - y_n).
    \end{align*}
}


\begin{definicija}
    \textbf{Skalarni proizvod} na \( \mathbb{R}^n \) definišemo sa \( <,> := \mathbb{R}^n
    \times \mathbb{R}^n \rightarrow \mathbb{R}\), odnosno

    \[
        <(x_1, \dots, x_n), (y_1, \dots, y_n)> = x_1y_1 + \dots + x_ny_n
         := \sum_{i = 1}^{n}{x_iy_i}.
    \]
\end{definicija}

\begin{definicija}
    \textbf{Norma vektora} \( x = (x_1, \dots, x_n) \in \mathbb{R}^n \) je 
    \[
        \lVert x \rVert = \lVert (x_1, \dots, x_n) \rVert := \sqrt{<x,x>} 
        = \sum_{i=1}^{n}{x_i^2}
    \]
\end{definicija}

\begin{primer}
    \( x = (2, 1, -3), \; y = (0, 0, 1) \):
    \tdplotsetmaincoords{70}{105}
    \begin{center}
        \begin{tikzpicture}[tdplot_main_coords, scale=1]
            
            % Ose
            \draw[->] (0,0,0) -- (3,0,0) node[anchor=north east]{$x_1$};
            \draw[->] (0,0,0) -- (0,2,0) node[anchor=north west]{$x_2$};
            \draw[->] (0,0,-4) -- (0,0,2) node[anchor=south]{$x_3$};

            % Vektori
            \draw[->, thick, blue] (0,0,0) -- (2,1,-3) node[below right] {$x$};
            \draw[->, thick, red] (0,0,0) -- (0,0,1) node[right] {$y$};

            % Krajnje tačke
            \fill (2,1,-3) circle (0.8pt) node[below left] {$(2,1,-3)$};
            \fill (0,0,1) circle (0.8pt) node[below right] {$(0,0,1)$};

            % Projekcije za x
            \draw[dashed] (2,1,-3) -- (2,1,0);
            \draw[dashed] (2,1,0) -- (2,0,0);
            \draw[dashed] (2,1,0) -- (0,1,0);

        \end{tikzpicture}
    \end{center}

    Ovde je skalarni proizvod \( <x, y> = -3 \);
    
    Norma vektora \( x \) je  \( \lVert x \rVert= \sqrt{2^2 + 1^2 + 3^2} 
    = \sqrt{4 + 1 + 9} = \sqrt{14} \);

    Norma vektora \( y \) je \( \lVert y \rVert= \sqrt{0^2 + 0^2 + 1^2} 
    = \sqrt{0 + 0 + 1} = \sqrt{1} = 1 \);

\end{primer}


\begin{definicija}
     Norme definišu \textbf{rastojanje između dve tačke} sa oznakom \(d(x, y) = 
     \lVert x-y \rVert \) kao
     \[
        \lVert x - y \rVert := \sqrt{(x_1 - y_1)^2 + \dots + (x_n - y_n)^2} 
        = \sum_{i=1}^{n}{(x_i - y_i)^2}.
     \] 
\end{definicija}

\subsubsection{Neka svojstva skupova u \( \mathbb{R}^n \)}

\begin{definicija}
    \textbf{Otvorena kugla} ili \textbf{otvorena lopta} u \(\mathbb{R}^n\) sa centrom u \(a\), poluprečnikom
    \(r > 0\) je skup
    \[
        B(a,r) := \{ x \in \mathbb{R}^n \mid \lVert x - a \rVert < r \}.
    \]
\end{definicija}

\begin{primer}
    U \( R^2 \), za tačku \( a = (2, 3) \) i poluprečnik \( r = 1 \), otvorena kugla
    izgleda ovako: 
    \begin{center}
        \begin{tikzpicture}[scale=0.91]
            % Ose
            \draw[->] (0,0) -- (5,0) node[right] {$x$};
            \draw[->] (0,0) -- (0,5) node[above] {$y$};

            % Mreža (opciono)
            \draw[step=1, gray!30, very thin] (0,0) grid (5,5);

            % Otvorena kugla
            \fill[blue!15] (2,3) circle (1);
            \draw[blue, thin, dashed] (2,3) circle (1);

            % Centar
            \fill (2,3) circle (1.2pt);
            \node[above right] at (2,3) {$a=(2,3)$};

            % Poluprečnik
            \draw[red, thick] (2,3) -- (3,3);
            \node[below] at (2.5,3) {$r=1$};
        \end{tikzpicture}
    \end{center}
    Tačka \( (x,y)\) pripada ovoj otvorenoj kugli samo ako važi \( \sqrt{(x-2)^2+(y-3)^2} < 1 \).
\end{primer}


\begin{definicija}
    \textbf{Zatvorena kugla} ili \textbf{zatvorena lopta} u \(\mathbb{R}^n\) sa centrom u \(a\),
    poluprečnikom \(r > 0\) je skup
    \[
        B[a,r] := \{ x \in \mathbb{R}^n \mid \lVert x - a \rVert \leq r \}.
    \]
\end{definicija}


\begin{primer}
    U \( \mathbb{R}^2 \), za tačku \( a = (1, 1) \) i poluprečnik \( r = 2 \), zatvorena kugla
    izgleda ovako:

    \begin{center}
        \begin{tikzpicture}[scale=0.91]

            % Ose
            \draw[->] (-2,0) -- (5,0) node[right] {$x$};
            \draw[->] (0,-2) -- (0,5) node[above] {$y$};

            % Mreža
            \draw[step=1, gray!30, very thin] (-2,-2) grid (5,5);

            % Zatvorena kugla
            \fill[blue!15] (1,1) circle (2);
            \draw[blue, thick] (1,1) circle (2);

            % Centar
            \fill (1,1) circle (1.2pt);
            \node[above right] at (1,1) {$a=(1,1)$};

            % Poluprečnik
            \draw[red, thick] (1,1) -- (3,1);
            \node[below] at (2,1) {$r=2$};

        \end{tikzpicture}
    \end{center}
    Zatvorena kugla sadrži i svoju ivicu za razliku od otvorene kugle. Tačka \( (x,y)\) joj pripada
    samo ako važi \( \sqrt{(x-1)^2 + (y-1)^2} \leq 2 \).
\end{primer}

\begin{definicija}
    \textbf{Rub otvorene i zatvorene kugle} u \(\mathbb{R}^n\) je 
    \[
        \partial B(a,r) = \partial B[a,r] = \{ x \in \mathbb{R}^n \mid \lVert x - a \rVert = r \}
    \]
\end{definicija}

\begin{definicija}
    Neka je skup \(A \subseteq \mathbb{R}^n\). Za tačku \(a \in \mathbb{R}^n\) kažemo
    da je \textbf{rubna tačka} skupa \(A \) ako za svako \(r > 0\) važi
    \begin{align*}
        B(a, r) \cap A &\neq \emptyset \\
        B(a, r) \cap A^c &\neq \emptyset
    \end{align*}
\end{definicija}

\begin{primer}
    Ovako izgleda definicija odozgo:
    \begin{center}
        \begin{tikzpicture}[scale=1.1]

            % Osa (opciono)
            \draw[->, gray] (-0.5,0) -- (7,0);
            \draw[->, gray] (0,-0.5) -- (0,5.5);

            % Skup A
            \fill[red!15]
                (1.2,1.3)
                .. controls (1.8,3.8) and (3.2,4.3) ..
                (4.8,3.8)
                .. controls (6.0,3.3) and (5.8,1.8) ..
                (4.9,1.3)
                .. controls (3.8,0.8) and (2.1,0.9) ..
                (1.2,1.3);
            \draw[red, thick]
                (1.2,1.3)
                .. controls (1.8,3.8) and (3.2,4.3) ..
                (4.8,3.8)
                .. controls (6.0,3.3) and (5.8,1.8) ..
                (4.9,1.3)
                .. controls (3.8,0.8) and (2.1,0.9) ..
                (1.2,1.3);

            % Oznaka skupa A
            \node[red!70!black] at (3.3,2.5) {$A$};

            % Rubna tacka a
            \coordinate (a) at (4.75,3.72);
            \fill[black] (a) circle (1.2pt);
            \node[above right] at (a) {$a$};

            % Kugle oko tacke a
            \draw[blue, thick] (a) circle (0.9);
            \draw[blue, thick] (a) circle (0.55);
            \draw[blue, thick] (a) circle (0.28);

            % Oznaka jedne kugle
            \node[blue] at (5.9,4.45) {$B(a,r)$};

        \end{tikzpicture}
    \end{center}

    Na slici je crvenom bojom predstavljen skup \(A\), dok je tačka \(a\) rubna tačka skupa \(A\).
    Za svako \(r>0\), otvorena kugla \(B(a,r)\) sadrži i tačke koje pripadaju skupu \(A\)
    i tačke koje ne pripadaju skupu \(A\), odnosno važi
    \[
        B(a,r)\cap A \neq \emptyset
        \qquad \text{i} \qquad
        B(a,r)\cap A^c \neq \emptyset.
    \]
\end{primer}

\begin{definicija}
    \textbf{Granica (rub)} skupa \(A \subseteq\) je skup svih rubnih tačaka s oznakom
    \(\partial A\).
\end{definicija}

\begin{definicija}
    Skup \(A \subseteq \mathbb{R}^n \) je \textbf{otvoren skup} ako važi
    \[
        (\forall a \in A) (\exists r > 0) : \ B(a,r) \subset A,
    \]
    odnosno da za svaku tačku u skupu \(A\) možemo napraviti kuglu oko nje koja je cela unutar
    skupa \(A\).
\end{definicija}

\begin{definicija}
    Skup \(A \subseteq \mathbb{R}^n \) je \textbf{zatvoren skup} ako je \(A^c\) otvoren skup.

    Skup \(A \subseteq \mathbb{R}^n \) je takođe zatvoren skup ako sadrži sve svoje 
    rubne tačke, odnosno \( \partial A \subseteq A \).
\end{definicija}

\subsection{Nizovi u \(\mathbb{R}^k\)}

Pošto \(n\) koristimo za indeks niza \(a_n\), \(k\) će nam označavati dimenziju prostora.

\begin{definicija}
    Niz \(a_n \in \mathbb{R}^k, n \in \mathbb{N} \) je niz uređenih \(k\)-torki, \(k \geq 1\)
    \[
        a_n = (\overset{1}{a_n}, \overset{2}{a_n}, \dots, \overset{k}{a_n}).
    \]
\end{definicija}

\begin{primer}
    \(k = 2\):
    \[
        a_n = (\frac{1}{n}, n) \in \mathbb{R}^2.
    \]
\end{primer}

\begin{definicija}
    Niz \(a_n\) u \(\mathbb{R}^k\) \textit{konvergira} ka \(a \in \mathbb{R}^k \) ako važi
    \[
        (\forall \varepsilon > 0) (\exists n_0 \in \mathbb{N}) 
        (\forall n \geq n_0 \implies \lVert a_n - a \rVert < \varepsilon) 
    \]
    \begin{center}
        odnosno drugačijim zapisom 
    \end{center}
    \[
        (\forall \varepsilon > 0) (\exists n_0 \in \mathbb{N}) 
        (\forall n \geq n_0 \implies d(a_n, a) < \varepsilon) 
    \]
    i ima oznaku
    \[
        \lim_{n \to \infty} {a_n} = a.
    \]
\end{definicija}
    
\begin{definicija}
    Ako niz ne konvergira, kažemo da \textit{divergira}.
\end{definicija}

\begin{stav}
    Za niz \(a_n = (\overset{1}{a_n}, \overset{2}{a_n}, \dots, \overset{k}{a_n}) \in \mathbb{R}^k\)
    i \(a = (a_1, a_2, \dots, a_k) \in \mathbb{R}^k\) važi
    \[
        \lim_{n \to \infty} a_n = a \in \mathbb{R}^k \iff 
        \lim_{n \to \infty} \overset{1}{a_n} = a_1 \dots,
        \lim_{n \to \infty} \overset{k}{a_n} = a_k
    \]
\end{stav}
\begin{dokaz}
    $\\$
    \(\implies\)
    
    Važi \(\lim_{n \to \infty} a_n = a \in \mathbb{R}^k\). 


    Neka je \(i \in {1, \dots, k}\) proizvoljni indeks. Pokažimo
    \(\lim_{n \to \infty} \overset{i}{a_n} = a_i.\)


    Neka je \(\varepsilon > 0\) proizvoljno. Tada 
    \((\forall \varepsilon > 0)(\exists n_0 \in \mathbb{N})
        (\forall n \geq n_0 \implies \lVert a_n - a \rVert < \varepsilon )\)
    \begin{align*}
        \lvert \overset{i}{a_n} - a_i \rvert &= \sqrt{(\overset{i}{a_n} - a_i)^2} \\
        &\leq \sqrt{(\overset{1}{a_n} - a_1)^2 + \dots +
        (\overset{k}{a_n} - a_k)^2} \\
        &= \lVert a_n - a \rVert .
    \end{align*}
    Iz pretpostavke konvergencije u Definiciji \textbf{1.19} važi da
    \(\lVert a_n - a \rVert < \varepsilon\). Tako da onda važi
    \[
        \lvert \overset{i}{a_n} - a_i \rvert \leq \lVert a_n - a \rVert  < \varepsilon.
    \]
    Dakle za isto \(n_0\) važi \(\forall n \geq n_0\).
    \[
        \lvert \overset{i}{a_n} - a_i \rvert < \varepsilon.
    \]
    $\\$
    \(\impliedby\)

    Pretpostavićemo da važi
    \[
        \lim_{n \to \infty} \overset{1}{a_n} = a_1 \dots,
        \lim_{n \to \infty} \overset{k}{a_n} = a_k
    \].
    Neka je \( \varepsilon > 0 \) proizvoljno. Treba naći \(n_0 \in \mathbb{N}\)
    takvo da \(\forall n \geq n_0\)
    \[
        \lVert a_n - a \rVert  < \varepsilon.
    \]
    Pošto smo pretpostavili da važe \(\lim_{n \to \infty} \overset{1}{a_n} = a_1, \dots,
    \lim_{n \to \infty} \overset{k}{a_n} = a_k\),
    onda
    \begin{align*}
    (\exists n_1 &\in \mathbb{N}) (\forall n \geq n_1)
    \lvert \overset{1}{a_n} - a_1 \rvert < \frac{\varepsilon}{k}\\
    &\vdots \\
    (\exists n_k &\in \mathbb{N}) (\forall n \geq n_k) 
    \lvert \overset{k}{a_n} - a_k \rvert < \frac{\varepsilon}{k}\\
    \end{align*}
    Neka sada uzmemo najveće \(n\), odnosno
    \[
        n_0 = \max\{n_1, \dots, n_k\}.
    \]
    Tada će za takvo \(n_0\) važiti i sve prethodne \(k\) nejednakosti.

    \vspace{1em}
    Imamo sledeću nejednakost
    \[
        \sqrt{(\overset{1}{a_n} - a_1)^2 + \dots +
        (\overset{k}{a_n} - a_k)^2}
        \leq
        \lvert \overset{1}{a_n} - a_1 \rvert + \dots +
        \lvert \overset{k}{a_n} - a_k \rvert \tag{1}
    \].
    Desni deo nejednakosti je veći kad bismo kvadrirali pa možemo
    ovako postaviti (svi su nenegativni pa smemo).

    Sada za svako \(n\) veće od izabranog maksimuma \(n_0\), tojest
    \(\forall n \geq n_0\)
    \begin{align*}
        \sqrt{(\overset{1}{a_n} - a_1)^2 + \dots +
        (\overset{k}{a_n} - a_k)^2}
        &\leq
        \lvert \overset{1}{a_n} - a_1 \rvert + \dots +
        \lvert \overset{k}{a_n} - a_k \rvert. \\
        &< \frac{\varepsilon}{k} + \dots + \frac{\varepsilon}{k} \\
        &= \varepsilon.
    \end{align*}
\end{dokaz}

\napomena{ Kod nejednakosti \textbf{(1)} ideja je da dokažemo da je 
    \(\lVert a_n - a \rVert < \varepsilon\) koristeći proizvoljni 
    \(\varepsilon > 0\) i \(n_0\) koje imamo. Znamo da je \(\lVert a_n - a \rVert =
    \sqrt{(\overset{1}{a_n} - a_1)^2 + \dots +
        (\overset{k}{a_n} - a_k)^2} \) i da bismo došli do  \(\varepsilon\), sa druge
        strane koristimo  
        \(\lvert \overset{1}{a_n} - a_1 \rvert 
        \dots \lvert \overset{k}{a_n} - a_k \rvert.\) 
    }

\subsection{Funkcije u \(\mathbb{R}^k\)}

\begin{definicija}
    \(D_f\) je domen funkcije \(f: D_f \to \mathbb{R}\) (skup tačaka u kojima je 
    funkcija \(f\) definisana).
\end{definicija}

\begin{primer}
    \(f(x, y) = \sqrt{x + y}\)

    \(D_f = \{ (x + y) \in \mathbb{R}^2 \mid x + y \geq 0\}\).

    \begin{center}
        \begin{tikzpicture}
            \begin{axis}[
                xlabel={$x$},
                ylabel={$y$},
                axis lines=middle,
                xmin=-4, xmax=4,
                ymin=-4, ymax=4,
                xtick distance=1,
                ytick distance=1,
                grid=both,
                axis equal
            ]

            % Osenčena oblast domena: y >= -x
            \addplot[
                draw=none,
                fill=red!30!white,
                fill opacity=0.2
            ] coordinates {(-4,4) (4,4) (4,-4) (-4,4)};

            % Granica domena
            \addplot[red, thick, dashed, domain=-4:4] {-x};

            \node at (axis cs:2,3) {$x+y\ge 0$};
            \node at (axis cs:2.2,-1.8) {$y=-x$};

            \end{axis}
        \end{tikzpicture}
    \end{center}

\end{primer}

\begin{definicija}
    \(a \in \mathbb{R}^k\) je \textbf{tačka nagomilavanja} skupa 
    \(A \subseteq \mathbb{R}^k\) ako važi
    \[
        (\forall \varepsilon > 0) (\exists B(a, \varepsilon)) 
        A \cap (\underbrace{B(a, \varepsilon) \setminus 
        \{a\}}_{\text{probušena okolina}}) \neq \emptyset
    \]
\end{definicija}

\begin{definicija}
    Neka je data funkcija \(f: D_f \to \mathbb{R}\) (\(D_f \subseteq \mathbb{R}^k\)) i
    neka je \(a \in \mathbb{R}^k\) tačka nagomilavanja skupa \(D_f\), tada važi
    \[
        \lim_{x \to a} f(x) = L \in \mathbb{R}.
    \]
    Tojest funkcija \(f\) konvergira ka \(L \in \mathbb{R}\) kada \(x \to a\)
    ako važi
    \[
        (\forall \varepsilon > 0) (\exists \delta > 0) (\forall x \in D_f)
        (\underbrace{x \in B(a, \delta) \setminus \{a\}}_
        {0 < \lVert x - a \rVert < \delta} \implies \lvert f(x) - L \rvert < \varepsilon).
    \]
    Kažemo da je \(L \in \mathbb{R}\) \textbf{limes funkcije} \(f\) kada \(x\) teži \(a\).
\end{definicija}

\napomena{tačka \(a\)
    ne mora da pripada \(D_f\).}


\begin{stav}
    Funkcije \(f: D_f \to \mathbb{R}\) (\(D_f \subseteq \mathbb{R}^k\)) i
    \(g: D_g \to \mathbb{R}\) (\(D_g \subseteq \mathbb{R}^k\)), neka je 
    \(a \in \mathbb{R}^k\) tačka nagomilavanja skupova \(D_f\) i \(D_g\). Neka je
    \(\lim_{x \to a}f(x) = L \in \mathbb{R} \text{ i } 
    \lim_{x \to a}g(x) = S \in \mathbb{R}\). Tada važi
    \begin{align}
        &\lim_{x \to a} (f\pm g)(x) = L + S, \\
        &\lim_{x \to a} (f \cdot g)(x) = L \cdot S, \\
        &\lim_{x \to a} \frac{f(x)}{g(x)} = \frac{L}{S} 
        \text{, ako \(\exists B(a, r) (\forall x \in B(a,r) \cap D_g)\)
         \(g(x) \neq 0\)}.
    \end{align}
\end{stav}

\begin{stav}
    \textbf{(Hajneov princip)}
    Neka je \(f: D_f \to \mathbb{R}\), \(a \in \mathbb{R}^k\) tačka nagomilavanja
    skupa \(D_f\). Tada važi
    \[
        \lim_{x \to a} f(x) = L \iff \forall a \in \mathbb{R}^k 
        \lim_{n \to \infty} a_n = a \implies \lim_{n \to \infty} f(a_n) = L. 
    \]
\end{stav}


\subsection{Neprekidnost funkcije}

\begin{definicija}
    Neka je \(D_f \to \mathbb{R}\) (\(D_f \subseteq \mathbb{R}^k\)). Neka je 
    \(a \in D_f\) tačka nagomilavanja skupa \(D_f\). Kažemo da je
    funkcija \(f\) neprikidna u tački \(a \in D_f\) ako važi
    \[
        \lim_{x \to a}f(x) = f(a).
    \]
\end{definicija}

\begin{definicija}
    Kažemo da je funkcija \(f\) neprekidna na svom domenu ako je neprekidna
    u svakoj tački nagomilavanja.
\end{definicija}

\begin{stav}
    Ako su funkcije \(f: D_f \to \mathbb{R}\) i \(g: D_g \to \mathbb{R}\) neprekidne
    u tački \(a \in D_f \cap D_g\), tada je funkcija
    \[
        \text{ neprekidna u } a
        \left\{
        \begin{aligned}
        f \pm g, \\
        f \cdot g, \\
        \frac{f}{g} \quad
        \text{ako } \exists B(a,r) \text{ td. }
        \forall x \in B(a,r) \ g(x) \neq 0.
        \end{aligned}
        \right.
    \]
\end{stav}


\begin{stav}
    \(f_1: D \to \mathbb{R}, f_2: D \to \mathbb{R}, \dots , f_k: D \to \mathbb{R}\)

    \(f:D \to \mathbb{R}^k\), \(f(x) = (f_1(x), f_2(x), \dots , f_k(x))\)

    \(g: \mathbb{R}^k \to \mathbb{R}\).
    
    Neka je \(f_1, \dots, f_k\) neprekidna u \(a \in D\), neka je \(b = (f_1(a), \dots, f_k(a))\)
    
    Neka je funkcija \(g\) neprekidna u \(b\), 
    tada je funkcija \(g \circ f: \mathbb{R}^k \to \mathbb{R}\) neprekidna u tački \(a\).
\end{stav}

\begin{dokaz}
    \(\\\)
    Neka je \(a_n \in \mathbb{R}^k\) proizvoljan niz takav da 
    \(\lim_{n \to \infty}a_n = a\).
    
    Dokazujemo \(\lim_{n \to \infty}g(f(a_n)) = g(f(a))\).

    Funkcije \(f_1, \dots, f_k\) su neprekidne u \(a\).
    \begin{align*}
        &\overset{\text{Hajne}}{\implies} \lim_{n \to \infty} f_1(a_n) = f_1(a), \dots ,
        \lim_{n \to \infty} f_k(a_n) = f_k(a)\\
        &\overset{\text{stav 2. čas}}{\implies} \lim_{n \to \infty}(f_1(a_n), 
        \dots, f_k(a_n)) = (f_1(a), \dots, f_k(a)) \\
        &\implies \lim_{n \to \infty} f(a_n) = f(a) = b \\
        &\overset{\text{g neprekidna i Hajne}}{\implies} \lim_{n \to \infty} g(f(a_n)) = g(b).
    \end{align*}

\end{dokaz}

\subsection{Vektorske funkcije}

\begin{definicija}
    \textbf{Vektorska funkcija} je funkcija \(F:D_F \to \mathbb{R}^k\) 
    \((D_f \subseteq \mathbb{R}^n)\). 
\end{definicija}

\begin{primer}
    \(F(x, y) = (x^2 + y, \sin{x}), \quad F: \mathbb{R}^2 \to \mathbb{R}^2 \).
\end{primer}

\begin{definicija}
    Neka je \(F: D_f \to \mathbb{R}^k\) (\(D_f \subseteq \mathbb{R}^n\))
    vektorska funkcija i \(a \in \mathbb{R}^n\) tačka nagomilavanja. Tada
    je \(F(x) = L \in \mathbb{R}^k\) ako i samo ako

    \[
        (\forall \varepsilon > 0) (\exists \delta > 0) 
        (\forall x \in B(a, \delta) \setminus \{a\} \cap D_f) 
        f(x) \in B(L, \varepsilon)
    \]
    tojest drugačije
    \[
        (\forall \varepsilon > 0) (\exists \delta > 0)
        (\forall x \in D_f) 0 < \lVert x - a \rVert < \delta
        \implies \lVert f(x) - L \rVert < \varepsilon.
    \]
\end{definicija}

\begin{definicija}
    
    Funkcija \(F: D_f \to \mathbb{R}^k\) (\(D_f \subseteq \mathbb{R}^n\))
    je neprekidna funkcija u \(a \in D_f \) tačka nagomilavanja ako važi
    \[
        \lim_{x \to a}F(x) = F(a).
    \]
\end{definicija}

\begin{primer}
    Naći \(D_f\):

    \begin{align*}
        F(x, y) &= (ln(x-y), \sqrt[3]{\cos x}) \\
        x - y &> 0 \\
        x &> y \\
        D_f &= \{(x,y) \in \mathbb{R}^2 \mid x > y \} \text{ - otvoren skup.}
    \end{align*}
\end{primer}


\subsection{Parcijalni izvodi}

\begin{definicija}
    Prvi parcijalni izvod funkcije \(f\) u \(x = (x_0, y_0)\) je
    \[
        \frac{\partial f}{\partial x_k} (x) \overset{\text{def}}{=}
        \lim_{h \to 0} \frac{f(x_0 + h, y_0) - f(x_0, y_0)}{h}. 
    \]
\end{definicija}


\subsection{Izvod u pravcu}

Neka je \(f: D_f \to \mathbb{R}\) (\(D_f \subseteq \mathbb{R}^n\)) 
i neka je \(v \in \mathbb{R}^k\).

%TODO - primer

\begin{definicija}
    \textit{Izvod} funkcije \(f\) \textit{u pravcu vektora} \(v = (v_1, \dots, v_n)\in \mathbb{R}^n\)
    u tački \(x = (x_1, \dots, x_n) \in D_f\) je 
    \[
        f'_v(x) = \lim_{h \to 0} \frac{f(x+hv) - f(x)}{h}, \quad h\in \mathbb{R}.
    \]
    Dakle
    \[
        f'_v(x) = \lim_{h \to 0} \frac{f(x_1 + hv_1, \dots x_n + hv_n) 
        - f(x_1, \dots , x_n)}{h}, \quad h\in \mathbb{R}
    \]
\end{definicija}



\begin{stav}
    \begin{align*}
        \frac{\partial f}{\partial x_1} (x) &= f'_v(x), \quad v = (1, 0, \dots , 0) \\
        &\vdots \\
        \frac{\partial f}{\partial x_n} (x) &= f'_v(x), \quad v = (0, 0, \dots , 1) \\
    \end{align*}
\end{stav}

%TODO - primer

\subsection{Diferencijabilnost funkcije f}


\textbf{Analiza 1:}

Za funkciju \(f: (a, b) \to \mathbb{R}\), kažemo da je diferencijabilna u 
\(x \in (a, b) \) ako 
\begin{align*}
    \exists \lim_{h \to 0} \frac{f(x + h) - f(x)}{h} = f'(x) &\iff \lim_{h \to 0}
    \frac{f(x+h) - f(x) - f'(x) h}{h} = 0 \\
    &\iff f(x+h) - f(x) - f'(x)h = o(h), \quad h \to 0 \\
    &\iff f(x+h) = f(x) + f'(x)h + o(h), \quad h \to 0
\end{align*}

\begin{definicija}
    Neka je \(fL D_f \to \mathbb{R}\) (\(D_f \subseteq \mathbb{R}^n\)) otvoren skup.
    Tada je funkcija \(f\) \textit{diferencijabilna} u tački \(x \in D_f\) ako postoji
    linearno preslikavanje \(L: \mathbb{R}^n \to \mathbb{R}\) tako da važi
    \begin{align*}
        f(x+h) &= f(x) + L\cdot h + o(\lVert h \rVert), \quad h \to 0 \\
        &\text{tojest} \\
        f(x+h) &= f(x) + L_1h_1 + \dots + L_nh_n + o(\lVert h \rVert), \quad h \to 0 \\
        h \in \mathbb{R}^n, \quad h &= (h_1, \dots, h_n), \quad \lVert h \rVert = \sqrt{h^{2}_1 + \dots + h^{2}_n} \in \mathbb{R}.
    \end{align*}
\end{definicija}

%TODO - primer

\begin{teorema}
    Ako je \(f: D_f \to \mathbb{R}\) (\(D_f \subseteq \mathbb{R}^n\)) diferencijabilna
    u tački \(x \in D_f\), tada postoje 
    \[
        \frac{\partial f}{\partial x_1}(x), \dots, \frac{\partial f}{\partial x_n}(x)
    \]
    i važi 
    \[
        f(x+h) = f(x) + \frac{\partial f}{\partial x_1}(x)h_1 + \dots + 
        \frac{\partial f}{\partial x_n}(x)h_n + o(\lVert h \rVert), \quad h \to 0.
    \]
\end{teorema}

%TODO - dokaz

\begin{stav}
    \(f:D_f \to \mathbb{R}\) diferencijabilna u tački \(a \in D_f \implies f\)
    neprekidna u \(a \in D_f\).

\end{stav}
\napomena{Obrnuti smer ne važi, odnosno ako je \(f: D_f \to \mathbb{R}\) neprekidna funkcija,
ne mora da znači da je i diferencijabilna.}

\napomena{Ako funkcija ima sve parcijalne izvode, ne mora ni tada značiti da je diferencijabilna.}
%TODO - dokaz

%TODO - primer

\begin{definicija}
    Neka je \(f: D_f \to \mathbb{R}\) (\(D_f \subseteq \mathbb{R}^n\)) funkcija koja
    ima sve parcijalne izvode 
    \(\frac{\partial f}{\partial x_1}, \dots, \frac{\partial f}{\partial x_n}\),
    \textit{gradijent} funkcije \(f\) u tački \(x \in D_f\) je vektor
    \[
        \triangledown f(x) = 
        \left(\frac{\partial f}{\partial x_1}, \dots, \frac{\partial f}{\partial x_n}\right) 
        \in \mathbb{R}^n.
    \]
\end{definicija}

\(\triangledown f\) je primer vektorskog polja.

%TODO - primer

\begin{teorema}
    Neka je funkcija \(f: D_f \to \mathbb{R}\) (\(D_f \subseteq \mathbb{R}^n\))
    diferencijabilna za svaku tačku domena, odnosno 
    \(\forall x \in D_f\). Tada \(\exists f'_v(x)\) za svaki vektor 
    \(v \in \mathbb{R}^n\)
    \begin{align*}
        f'_v(x) &= <\triangledown f(x), v>, \quad v \in \mathbb{R}^n,\, v = (v_1, \dots , v_n) \\
        f'_v(x) &= \frac{\partial f}{\partial x_1}(x)v_1 + \dots
        + \frac{\partial f}{\partial x_n}(x)v_n. 
    \end{align*}
\end{teorema}

\napomena{Obrnuto ne mora da važi!}

%TODO - primer

\begin{teorema}
    Neka je data funkcija \(f: D_f \to \mathbb{R}\) (\(D_f \subseteq \mathbb{R}^n\))
    takva da postoje svi parcijalni izvodi 
    \(\frac{\partial f}{\partial x_1}, \dots, \frac{\partial f}{\partial x_n}\)
    za \(\forall x \in D_f\) i neka su oni neprekidne funkcije. Tada je 
    \(f\) diferencijabilna u svakoj tački domena, odnosno \(\forall x \in D_f\).
\end{teorema}


\subsection{Diferencijabilnost vektorske funkcije}

\begin{definicija}
    Funkcija \(F: D_F \to \mathbb{R}^m\) (\(D_F \subseteq \mathbb{R}^n\)) je diferencijabilna
    u tački \(a = (a_1, \dots, a_n) \in D_F\) ako postoji linearno preslikavanje
    \(L: \mathbb{R}^n \to \mathbb{R}^m\) takvo da važi
    \[
        F(a+h) = F(a) + L\cdot h + o(h), 
        \quad h \to 0, \;h = (h_1, \dots, h_n) \in \mathbb{R}^n
    \]
\end{definicija}

\begin{definicija}
    F je diferencijabilna funkcija na \(D_F\) ako je diferencijabilna u svakoj tački 
    domena \(D_F\).
\end{definicija}

\begin{primer}
    \(\\\)
    \begin{align*}
        F(x,y) &= (2x + 3y, 4x + 5y), \; h = (h_1, h_2) \\
        &(2(x + h_1) + 3(y + h_2), 4(x + h_1) + 5(y + h_2)) \\
        &(2x + 3y, 4x + 5y) + L \cdot (h_1, h_2) + o(h) \\
        &L(h_1, h_2) = (2h_1 + 3h_2, 4h_1 + 5h_2).
    \end{align*}
    \napomena{\(o(h)\) je nula jer je funkcija \(F\) linearna.}
\end{primer}

\begin{teorema}
    Neka je funkcija \(F: D_F \to \mathbb{R}^m\) \(D_F \subseteq \mathbb{R}^n\) takva da
    \(F(x) = (f_1(x), \dots, f_m(x))\), \(f_1, \dots, f_m : D_F \to \mathbb{R}\), tada je
    \(F\) diferencijabilna u tački \(a \in D_F\) ako i samo ako su funkcije 
    \(f_1, \dots, f_m\) diferencijabilne u tački \(a \in D_f\) i važi
    \[
        L =
        \begin{pmatrix}
        \frac{\partial f_1}{\partial x_1} & \frac{\partial f_1}{\partial x_2} & \cdots & \frac{\partial f_1}{\partial x_n} \\
        \vdots & \vdots & \ddots & \vdots \\
        \frac{\partial f_m}{\partial x_1} & \frac{\partial f_m}{\partial x_2} & \cdots & \frac{\partial f_m}{\partial x_n}
        \end{pmatrix} 
    \]
\end{teorema}

\begin{dokaz}
    \(\\\)
    \(F\) je diferencijabilna u \(a \in D_f \):
    \begin{align*}
        &\iff F(a + h) = F(a) + L \cdot h + o(h), 
        \quad h \to 0, \; h \in \mathbb{R}^n \\
        &\iff \begin{pmatrix}
                f_1(a + h) \\
                \vdots \\
                f_m(a + h)
            \end{pmatrix}
        =
            \begin{pmatrix}
                f_1(a) \\
                \vdots \\
                f_m(a)
            \end{pmatrix}
        +
            \begin{pmatrix}
                a_{11} & \cdots & a_{1n} \\
                \vdots & \ddots & \vdots \\
                a_{m1} & \cdots & a_{mn}
            \end{pmatrix}
            \begin{pmatrix}
                h_1 \\
                \vdots \\
                h_m
            \end{pmatrix}
        +
            \begin{pmatrix}
                o(\|h\|) \\
                \vdots \\
                o(\|h\|)
            \end{pmatrix} \\
        &\iff f_j(a+h) = f_j(a) + \sum_{i = 1}^{n}{a_{ji}\cdot h_i} + o(\lVert h \rVert),
        \quad h \to 0, \; \forall j = 1, \dots, m \\
        &\iff f_j \text{ je diferencijabilna } \forall j = 1, \dots, m.
    \end{align*} 

    Važi
    \[
        a_{ji} \text{ je } \frac{\partial f_j}{\partial x_i} (a)
        \text{(2. nedelja)}, \quad i = 1, \dots, n; \; j = 1, \dots, m.
    \]
\end{dokaz}

\begin{primer}
    Za prošli primer \((2x + 3y, 4x + 5y)\)
    \[
        L =
        \begin{pmatrix}
            2 & 3\\
            4 & 5
        \end{pmatrix}.
    \]
\end{primer}

\begin{definicija}
    Ako je funkcija \(F: D_F \to \mathbb{R}^m\) (\(D_F \subseteq \mathbb{R}^n\))
    diferencijabilna u tački \(a \in D_F\), tada za linearno preslikavanje 
    \(L: \mathbb{R}^n \to \mathbb{R}^m\) koristimo oznaku
    \[
        dF(a)
    \]
    i nazivamo \textbf{diferencijal} funkcije \(f\).
\end{definicija}
\napomena{\(f: D_f \to \mathbb{R}\)} ako je diferencijabilna u \(a \in D_f\), linearno
preslikavanje \(df(a): \mathbb{R}^n \to \mathbb{R}\) nazivamo diferencijal funkcije \(f\).


\begin{definicija}
    Matrica \(J = [\frac{\partial f_j}{\partial x_i}]\) se naziva 
    \textbf{Jakobijeva matrica} vektorske funkcije \(F\).
\end{definicija}
\napomena{Specijalno ako je \(F:D_F \to \mathbb{R}^n \), (\(D_F \subseteq \mathbb{R}^n\)),
onda je matrica dimenzije \(n \times x\), 
tada se determinanta \(detJ\) naziva \textbf{Jakobijan}.}


\begin{primer}
    \(\\\)    
    \begin{align*}
        F(x, y) &= (2x + 3y, 4x + 5y) \\
        det(J) &= 5 \cdot 2 - 3 \cdot 4 = 10 - 12 = -2.
    \end{align*}


\end{primer}

\begin{stav}
    Ako su funckije \(F, G: D \to \mathbb{R}^m\) (\(D \subseteq \mathbb{R}^n\))
    diferencijabilne u \(a \in D\) tada je i \(\lambda_1 F + \lambda_2G\) diferencijabilna
    u tački \(a \in D\)
    \[
        d(\lambda_1F + \lambda_2G)(a) = \lambda_1 \cdot dF(a) + \lambda_2 \cdot dG(a).
    \]
\end{stav}

\begin{dokaz}
    \begin{align*}
        (\lambda_1F + \lambda_2G)(a + h) &= \lambda_1 \cdot F(a+h) + \lambda_2 \cdot G(a+h) \\
        &= \lambda_1 \cdot (F(a) + dF(a)\cdot h + o(h)) + 
        \lambda_2 \cdot (G(a) + dG(a) \cdot h + o(h)) \\
        &= (\lambda_1F + \lambda_2G)(a) + (\underbrace{\lambda_1dF(a) + \lambda_2dG(a)}
        _{d(\lambda_1F + \lambda_2G)(a)})\cdot h + 
        \underbrace{\lambda_1o(h) + \lambda_2o(h)}_{o(h)}, \quad h \to 0
    \end{align*}    
\end{dokaz}

\begin{stav}
    Ako su \(fg: D \to \mathbb{R}^n\) (\(D \subseteq \mathbb{R}^n\)) diferencijabilne
    u tački \(a \in D\) tada je i
    \begin{enumerate}
        \item \(f \cdot g\) diferencijabilna u \(a \in D\) i važi:
        \[
            d(fg)(a) = f(a) \cdot dg(a) + g(a) \cdot df(a),
        \]
        \item ako \(g(x) \neq 0\) \(\forall x \in D\) tada je i \(\frac{f}{g}\)
        diferencijabilna u \(a \in D\) i važi:
        \[
            d\left(\frac{f}{g}\right)(a) = \frac{1}{g^2(a)}(g(a)\cdot df(a) - f(a) \cdot dg(a)).
        \]
    \end{enumerate}
\end{stav}

\subsection{Diferencijabilnost kompozicije}


\begin{definicija}
    Neka je funkcija \(F:D_F \to \mathbb{R}^m\) (\(D \subseteq \mathbb{R}^n\)) diferencijabilna
    u tački \(a \in D_F\) i neka je funkcija \(G: D_G \to \mathbb{R}^p\) 
    (\(D_G \subseteq \mathbb{R}^m\)) diferencijabilna u \(F(a)\) i \(F(D_F) \subseteq D_G\),
    tada je kompozicija \(G \circ F: D_F \to \mathbb{R}^p \) diferencijabilna u tački
    \(a \in D_F\) i važi
    \[
        d(G \circ F)(a) = dG(F(a)) \cdot dF(a).
    \]
\end{definicija}
\napomena{\(\cdot\) je matrično množenje.}
\begin{dokaz}
    Treba pokazati \( G \circ F (a+h) = G\circ F(a) + 
    \underbrace{d(G\circ F)(a)}_{dG(f(a)) \cdot dF(a)} \cdot h + o(h),
    \quad h \to 0, \; h \in \real^n\) i treba pokazati
    \(d(G\circ F)(a) = dG(F(a)) \cdot dF(a).\)

    \begin{align*}
        G\circ F(a+h) - G\circ F(a) &= G(F(a+h)) - G(F(a)) \\
        \text{Neka je } k = F(a+h) - F(a), \quad k \in \real^m \\
        G(F(a) - k) - G(F(a))& \\
    \end{align*}
\end{dokaz}

\begin{primer}
    \(\\\)
    Neka su \(F: \real^2 \to \real^3, \; G: \real^3 \to \real \):
    \begin{align*}
        F(x, y) &= (x^2, 3x + 2y, 4y) \\
        G(z_1, z_2, z_3) 7 &= (z_1 \cdot z_2 + z_3) \\
    \end{align*}
    Na osnovu teoreme \(G \circ F\) je diferencijabilna \(\forall x, y \in \real^2\) 
    (zato što su i \(F\) i \(G\)). Nađimo \(d(G\circ F)(x,y)\):
    \begin{enumerate}[label = \Roman*]
        \item način
            \begin{align*}
                G\circ F (x,y) &= x^2 (3x + 2y) + 4x = 3x^3 + 2x^2y + 4x \\
                d(G\circ F)(x,y) &= \left[ 9x^2 + 4xy + 4, \; 2x^2 \right]
            \end{align*}
    
        \item način preko teoreme
            \begin{align*}
                d(G \circ F)(x,y) &= dG(F(x,y)) \cdot dF(x,y) \\
                dG(z_1, z_2, z_3) &= \left[ z_2, z_1, 1 \right] \\
                &= \left[ 3x+2y, x^2, 1 \right] \cdot
                \begin{bmatrix}
                    2x & 0 \\
                    3 & 2  \\
                    4 & 0
                \end{bmatrix} \\
                &= (3x+2y)\cdot 2x + 3x^2 +4, \; 2x^2 =
                \begin{bmatrix}
                    9x^2 + 4xy + 4 & 2x^2
                \end{bmatrix}.
            \end{align*}
    \end{enumerate}
\end{primer}

\subsection{Izvodi inverzne funkcije}

\begin{teorema}
    Neka je funkcija \(F: D \to \real ^n\) (\(D \subseteq \real^n\)) diferencijabilna
    i neka je funkcija \(F : D \to F(D)\) bijekcija i neka je
    \(F^-1: F(D) \to D\) diferencijabilna, tada je
    \[
        dF^{-1} (F(x)) = d(F(x))^{-1}.
    \]
\end{teorema}
\napomena{\(F: \real \to \real \) bijekcija i diferencijabilna, tada postoji
\(F: \real \to \real \) koja ne mora biti diferencijabilna.

Primer: \(f(x) = x^3\)

\(f^{-1}(x) = \sqrt[3]{x} - \) nije diferencijabilna u tački \(x = 0\).}
\begin{dokaz}
    \begin{align*}
        F \circ F^{-1} &= id: \real^n \to \real^n, \quad id(x_1, \dots, x_n) = (x_1, \dots, x_n) \\
        F^{-1} \circ F &= id \\
        d(F \circ F^{-1})(F(x)) &= dF(x) \cdot dF^{-1}(F(x)) =
        \begin{bmatrix}
            1 &  & 0 \\
            &\ddots & \\
            0 &  & 1 
        \end{bmatrix} \\
        d(F^{-1}\circ F)(x) &= dF^{-1}(F(x))dF(x) = 
        \begin{bmatrix}
            1 &  & 0 \\
            &\ddots & \\
            0 &  & 1 
        \end{bmatrix} \\
        dF^{-1}(F(x)) &= (dF(x))^{-1}.
    \end{align*}
\end{dokaz}

\begin{primer}
\(F(x,y) = (2x+3y,4x+5y)\)
\begin{enumerate}[label = \Roman*]
    \item način
    \begin{align*}
        (dF^{-1})(F(x)) = -\frac{1}{2} \cdot 
        \begin{bmatrix}
            5 & -3 \\ -4 & 2
        \end{bmatrix} =
        \begin{bmatrix}
            \frac{-5}{2} & \frac{3}{2} \\ 2 & -1   
        \end{bmatrix}
    \end{align*}
    \item način Domaći
\end{enumerate}
\end{primer}

\subsection{Tangenta i tangentna ravan}

\begin{definicija}
    Neka je r\(r \; I \to \real^3\) (\(I \subseteq \real\)) diferencijabilno preslikavanje,
    tada se \(r\) ili \(r(I)\) naziva \textbf{kriva} u \(\real^3\).
\end{definicija}

\begin{definicija}
    Ako je \(r: I \to \real^3\) neprekidno diferencijabilno preslikavanje, onda se naziva
    \textbf{glatka kriva}.
\end{definicija}


\begin{definicija}
    Ako je \(r'(t) = (x'(t), y'(t), z'(t)) \neq (0, 0, 0) \; \forall t \in I\), onda je 
    \(r\) \textbf{regularna kriva}.
\end{definicija}

\begin{definicija}
    Neka je \(t\) tangenta na reglarnu krivu \(r\) u tački \(t_o\) ako prolazi kroz
    \(r(t_0)\) i koeficijent pravca joj je \((x'(t), y'(t), z'(t))\).
\end{definicija}

\textbf{LAAG:}
\[
    \frac{x-x_0}{a}= \frac{y - y_0}{b} = \frac{z - z_0}{c}
\]
gde su \((a, b, c) \in \real^3\) koeficijenti pravca.

Dakle tangenta \(t\) je zadata sa
\[
    \frac{x-x(t_0)}{x'(t_0)}= \frac{y - y(t_0)}{y'(t_0)} = \frac{z - z(t_0)}{z'(t_0)}
\]

\begin{primer}
    
    Neka je \(r: (-1, 1) \to \real^3\) i \(r(t) = (0, t, \sqrt{1 - t^2})\).
    Odredimo tangentu u \(t = 0\).

    \[\frac{x - 0}{0} = \frac{y - 0}{1} = \frac{z - 1}{0}\]
    znači
    \begin{align*}
        1\cdot (z-1) &= 0\cdot(y-0)\\
        z - 1 &= 0 \\
        z &= 1
    \end{align*}

    \begin{center}
        \begin{tikzpicture}[scale=2, line cap=round, line join=round]

            % ose
            \draw[->] (-1.2,0) -- (1.2,0) node[below left] {$y$};
            \draw[->] (0,-0.2) -- (0,1.4) node[left] {$z$};
            \draw[->] (0,0) -- (-0.7,-0.5) node[below] {$x$};

            % ravnina x=0 (tj. yz-ravan, samo nagovestaj)
            \draw[dashed, gray] (-1,0) -- (-1,1.2);
            \draw[dashed, gray] (1,0) -- (1,1.2);
            \draw[dashed, gray] (-1,1.2) -- (1,1.2);

            % kriva r(t) = (0,t,sqrt(1-t^2))
            % u yz-ravni: y=t, z=sqrt(1-t^2)
            \draw[thick, domain=-1:1, samples=100, smooth, variable=\t]
                plot ({\t},{sqrt(1-\t*\t)});

            % tangenta z=1, x=0
            \draw[thick] (-1,1) -- (1,1);

            % tacka dodira
            \fill (0,1) circle (0.8pt);
            \node[above right] at (0,1) {$(0,0,1)$};

            % oznaka tangente
            \node[right] at (0.75,1.03) {$T$};

        \end{tikzpicture}
    \end{center}
    \begin{align*}
        z (t) &= \sqrt{1-t^2} \\
        z(t) &= -\frac{z}{\sqrt{1-t^2}}
    \end{align*}
\end{primer}

\begin{primer}
   FALI %TODO - primer aa
\end{primer}

\begin{definicija}
    Neka je funkcija \(f: D_f \to \real\) (\(D_f \subseteq \real^3\)) diferencijabilna,
    tada se skup \(S = \{(x,y,z) \in \real^3 \mid f(x,y,z) = 0\}\) naziva \textbf{površi}.
\end{definicija}

\begin{definicija}
    Ako je \(f\) neprikidno diferencijabilna funkcija tada je \(S\) \textbf{glatka površ}.
\end{definicija}

\begin{definicija}
    Ako je gradijent \(\triangledown f(x,y,z) \neq (0,0,0)\) u svakoj 
    \((x,y,z) \in S\), tada je \(S\) \textbf{regularna površ}.
\end{definicija}

\begin{primer}
    \(f(x,y,z) = x^2+y^2+z^2=1\)
    \begin{align*}
        S &= \{(x,y,z) \in \real^3 \mid x^2+y^2+z^2=1\} \\
        \triangledown f(x,y,z) &= (2x,2y,2z) \neq (0,0,0) \quad \forall (x,y,z) \in S.
    \end{align*}
\end{primer}

\begin{definicija}
    Tangentna ravan na regularnu površ \(S\) u tački \(x_0, y_0, z_0\)
    je ravan koja sadrži tangente svih regularnih krivih \(r: I\to S\)
     (\(S \subseteq \real^3\)) koje prolaze kroz tačku \(x_0, y_0, z_0\).
\end{definicija}

\textbf{LAAG:}

Jednačina ravni kroz tačku \((x_0, y_0, z_0)\) kojoj je vektor normale
\(n = (n_1, n_2, n_3)\) se analitički zapisuje 
\(n_1(x-x_0) + n_2(y-y_0)+n_3(z-z_0) = 0\).

\begin{stav}
    Tangentna ravan na površ \(S\) u tački \(x_0, y_0, z_0\) ima vektor normale
    \(n = \triangledown f (x_0, y_0, z_0)\), tojest tangentna ravan je skup svih 
    pravih koje su ortogonalne na \(n = \triangledown f (x_0, y_0, z_0)\).
\end{stav}

\begin{dokaz}
    
    \begin{align*}
        r: I \to S &= \{(x,y,z) \mid f(x,y,z) = 0\} \\
        f \circ r(t) &= 0\\
        f \circ R: I \to \real (I \subseteq \real) \quad df(r(t)) \cdot r'(t) &=
        \begin{bmatrix}
            \frac{\partial f}{\partial x} & \frac{\partial f}{\partial y} & \frac{\partial f}{\partial z}
        \end{bmatrix}
        \cdot r(t) \cdot \begin{bmatrix}
            x'(t) \\ y'(t) \\ z'(t)
        \end{bmatrix} \\
        &= <\triangledown f(r(t), r'(t))> \\
        &= 0 \\
        \triangledown f(r(t)) &\perp r'(t) 
    \end{align*}
\end{dokaz}

\begin{primer}
    Naći tangentnu ravan na \(S = \{x^2+y^2+z^2=1\}\) u tački \(0,0,1\)
    \begin{align*}
        (n_1, n_2, n_3) &= \triangledown f(x,y,z) \\
        &= (2x,2y,2z) \\
        \triangledown f(0,0,1) &= (0,0,2) \\
        \text{Tangentna ravan} \\
        0(x-o) + 0(y-0) + 2(z-1) &= 0 \\
        z &= 1 \\
    \end{align*}
    \begin{center}
        \begin{tikzpicture}[scale=2.2, line cap=round, line join=round]

    % koordinate za "3D" izgled
    \def\R{1}
    \def\PlaneY{1}
    \def\PlaneH{0.32}
    \def\PlaneW{1.15}
    \def\TiltX{0.55}

    % osa z
    \draw[->] (0,-1.25) -- (0,1.45) node[left] {$z$};

    % osa y
    \draw[->] (-1.35,0) -- (1.35,0) node[right] {$y$};

    % osa x
    \draw[->] (0,0) -- (-0.85,-0.55) node[below] {$x$};

    % sfera
    \draw[thick] (0,0) circle (\R);

    % "ekvator" sfere za 3D utisak
    \draw[dashed] (-1,0) arc (180:360:1 and 0.28);
    \draw        (-1,0) arc (180:0:1 and 0.28);

    % tangentna ravan z=1
    \coordinate (A) at (-\PlaneW,\PlaneY);
    \coordinate (B) at (\PlaneW,\PlaneY);
    \coordinate (C) at ({\PlaneW+\TiltX},{\PlaneY+\PlaneH});
    \coordinate (D) at ({-\PlaneW+\TiltX},{\PlaneY+\PlaneH});

    \fill[gray!5] (A) -- (B) -- (C) -- (D) -- cycle;
    \draw[thick] (A) -- (B) -- (C) -- (D) -- cycle;

    % tačka dodira
    \fill (0,1) circle (0.9pt);
    \node[right] at (0.03,1.03) {$(0,0,1)$};

    % normala u tački dodira
    \draw[->, thick] (0,1) -- (0,1.35);

    % oznaka ravni
    \node[right] at (1.18,1.02) {$z=1$};

\end{tikzpicture}
    \end{center}
\end{primer}

\end{document}


